The TDECode class has the following class variables -

{\bf static TDEVPtype[] ptypes = {TDEVPtype.INPUT, TDEVPtype.OUTPUT, TDEVPtype.INOUT}} -
an array to convert a numeric port type of 0, 1 or 2 from a TDEPORT TDE
into a TDEVPtype enum used for entries in {\bf ext\_port\_dirs}
(described below).

{\bf static ArrayList\verb'<'Object\verb'>' ext\_ports} - a list of external
ports, each member is a Net or else a port identifier string.

{\bf static ArrayList\verb'<'TDEVPtype\verb'>' ext\_port\_dirs} - a list of
port types, i.e. input, output or 3-state output.

{\bf static ArrayList\verb'<'String\verb'>' ext\_port\_pins} - a list of port
pins where these are FPGA pads.

{\bf static TreeMap\verb'<'String,Object\verb'>' portmap} is a map of ports where
the key is either the port identifier of else the port Net.

The following member variables contain information about the target FPGA.

\blist
\item   {\bf String} fname - name of the FPGA family

\item   {\bf int} lutwidth - number of inputs to CLB LUT

\item   {\bf int} dspregwidth - effective width of multiplier/DSP
                                         block when used as a register

\item   {\bf int} srladdrwidth - shift register address width

\item   {\bf int} addsubinputwidth - maximum adder/subtracter input width

\item   {\bf int} multinputwidth - maximum multiplier input width

\item   {\bf int} divinputwidth - maximum divider input width

\item   {\bf int} ramb\_ports - maximum number of block RAM ports

\item   {\bf boolean[]} ramb\_readable - which block RAM ports are readable

\item   {\bf boolean[]} ramb\_writable - which block RAM ports are writable

\item   {\bf int[]} ramb\_dwidths = new int[20] - block RAM data widths for address widths

\item   {\bf int} ramb\_addrmin[] - block RAM smallest address width for \verb'#' of ports

\item   {\bf int} ramb\_addrmax[] - block RAM largest address width for \verb'#' of ports

\item   {\bf int} ramc\_ports - maximum number of CLB RAM ports

\item   {\bf boolean[]} ramc\_readable - which CLB RAM ports are readable

\item   {\bf boolean[]} ramc\_writable - which CLB RAM ports are writable

\item   {\bf int[]} ramc\_dwidths = new int[20] - CLB RAM data widths for address widths

\item   {\bf int} ramc\_addrmin[] - CLB RAM and ROM smallest address width for \verb'#' of ports

\item   {\bf int} ramc\_addrmax[] - CLB RAM port max block address width for \verb'#' of ports

\item   {\bf int} romc\_addrmin[] - CLB ROM port min block address width for \verb'#' of ports

\item   {\bf int} romc\_addrmax[] - CLB ROM port max block address width for \verb'#' of ports

\item   {\bf int[]} queue\_buffer\_cram\_depths - allowed depths for CLB RAM queue buffers

\item   {\bf int} queue\_buffer\_min\_cram\_depth - smallest depth for CLB RAM queue buffers

\item   {\bf int} queue\_buffer\_max\_cram\_depth - largest depth for CLB RAM queue buffers

\item   {\bf int[]} queue\_buffer\_rram\_depths - allowed depths for block RAM queue buffers

\item   {\bf int} queue\_buffer\_min\_rram\_depth - smallest depth for block RAM queue buffers

\item   {\bf int} queue\_buffer\_max\_rram\_depth - largest depth for block RAM queue buffers

\item   {\bf boolean} defaultOutputDirect - is the output mode direct flag default?

\item   {\bf boolean} defaultStaticContinuous - is the static mode continuous flag default?

\item   {\bf boolean} deviceHasTBUF - does device have 3-state buffers?

\item   {\bf boolean} deviceAllowsQueue - does device allow queue mode?
                                         queues require memory

\item   {\bf boolean} deviceAllowsCMemory - does device allow cmemory mode?
                                         requires CLB RAM

\item   {\bf boolean} deviceAllowsRMemory - does device allow rmemory mode?
                                         requires block RAM

\item   {\bf int} ramc\_mux\_addr\_bits\_maxc - CLB RAM address bits beyond maximum
                                         the number of extra address bits for
                                         CLB RAM gained by MUXing

\item   {\bf boolean} deviceFDRSEpatch - change FDRSE to FDCE + logic for CPLD

\item   {\bf boolean} deviceCascadeGates - skip conversion of cascaded or wide gates

\item   {\bf boolean} deviceCPLDarith - use CPLD addition/subtraction logic
                                         i.e. simple logic instead of carry chain
\blend

These variables are assigned in the constructor for the FPGA class, XC5V etc
which extend class TDECode.



The following methods return FPGA parameters derived from the above class
variables -
\blist
\item   LUTWidth() - return the LUT width for the FPGA
\item   AddSubWidth() - return the maximum adder/subtracter input
        width for the FPGA (9999 for no explicit limit)
\item   DivWidth() - return the maximum divider input width
        for the FPGA (9999 for no explicit limit)
\item   SRLAddrWidth() - return the maximum SRL (shift register)
        input width for the FPGA (0 if there is no SRL element)
\item   cramPorts() - return the number of memory ports for
        combinatorial-output RAM
\item   rramPorts() - return the number of memory ports for
        registered-output RAM
\item   cramAwidthMax() - return the maximum address width for a
        combinatorial-output RAM
\item   rramAwidthMax() - return the maximum address width for a
        registered-output RAM
\item   rramDwidth() - return the data width associated with an address
        width for a registered-output RAM
\item   cramReadable() - determine if a combinatorial-output RAM port is
        readable
\item   cramWritable() - determine if a combinatorial-output RAM port is
        writable
\item   rramReadable() - determine if a registered-output RAM port is
        readable
\item   rramWritable() - determine if a registered-output RAM port is
        writable
\item   queueDepthUsingCram() - determine if a requested queue buffer
        depth using combinatorial output memory is allowed
\item   queueDepthUsingRram() - determine if a requested queue buffer
        depth using registered output memory is allowed
\item   minQueueDepth() - return the minimum asynchronous queue depth
\item   maxQueueDepth() - return the maximum queue depth
\item   outputDirect() - determine if the output mode 'direct' attribute
        should be default for this device
\item   defaultContinuous() - determine if the selectvalue mode
        'continuous' attribute should be default for this device
\item   hasTBUF() - determine if this device has internal 3-state
        buffers
\item   allowQueue() - determine if queue mode variables are allowed for
        this device
\item   allowCMemory() - determine if cmemory mode variables are allowed
        forthis device
\item   allowRMemory() - determine if rmemory mode variables are allowed
        forthis device
\item   needFDRSEpatch() - determine if this device requires gates to
        fully implement an FDRSE flip-flop
\item   needCascadeGates() - determine if this device must construct
        wide AND or OR gates by cascading, rather than using inbuit
        MUXes or a carry chain
\item   needCPLDarith() - determine if special arithmetic for CPLDs is
        needed; this means that addition and subtraction and magnitude
        comparisons will be done by gate logic rather than using carry
        chains
\blend

Abstract method outputExterns() writes all external ports to the EDIF netlist
file. This method is overridden in class XTDECODE for Xilinx.

The following methods optimise the netlist code and are described
elsewhere.
\blist
\item   optimiseGNDandVCC()
\item   optimiseINV()
\item   optimiseAND()
\item   optimiseOR()
\item   optimiseXOR()
\item   unNest()
\item   inputEliminate()
\item   optimiseLUT()
\item   optimiseFDRS()
\item   remapInputs()
\blend









The remainder of the methods are abstract and are overridden by
methods in xilinx.XTDECode. All except the last seven are methods
to convert TDEs to internal netlist constructs. The method names
are self explanatory -
\blist
\item   TDESELECT()
\item   TDETSELECT()
\item   TDEREG()
\item   TDEINV()
\item   TDEAND()
\item   TDEOR()
\item   TDEXOR()
\item   TDEDFF()
\item   TDEIBUF()
\item   TDEOBUF()
\item   TDESTART()
\item   TDEILOOP()
\item   TDEWHEN()
\item   TDEWHILE()
\item   TDEDOWHILE()
\item   TDERESYNC()
\item   TDESAMPLE()
\item   TDEEXECP()
\item   TDECONNECT()
\item   TDEOPERATOR()
\item   TDEWAIT()
\item   TDEDEL()
\item   TDEDIVERGE()
\item   TDEQUEUEBUFFER()
\item   TDEPRIORITY()
\item   TDECRAM()
\item   TDERRAM()
\item   TDECLOCK()
\item   TDECLOCKINFO()
\item   TDEELEMENT()
\item   TDEELEMENTCONNECT()
\item   TDEPORT()
\item   TDESPECIAL()
\blend

The last seven abstract methods are miscellaneous. Again they are
overridden by methods in xilinx.XTDECode.

TDEOptimise() applies family-specific optimisations to the TDE list.

optimise() carries out low-level optimisation of the internal netlist
structure.

convert() converts generic gate elements to appropriate elements
for the target device. Before this is called the netlist structure
might contain an AND gate with 10 inputs. Following conversion that
element would have been replaced perhaps by 3 4-input LUTs connected
to carry chain MUXCYs.

verify() examines all elements to ensure there are no errors or
inconsistencies. This method is not located in class netlist.Element
as the checks are specific to target device elements rather than
generic elements (there is however a method netlist.Net.verify()
to check nets and this is generic since nets are not specific to
the target device).

netFromPin() returns the identifier of the net attached to a pad
given the pad identifier.

EDIFFileNameExtension() returns the EDIF file extension.


Abstract method outputConstraints() writes constraints to the approprite file(s)
(for Xilinx this is the ISE NCF or Vivado XDC file). This method is overridden in
class XTDECODE for Xilinx.
